% Abstract – resumo em língua estrangeira (NBR 14724).
% RASCUNHO escrito a partir do resumo em português, não traduzido palavra por palavra.
% Revise antes da entrega (de preferência com alguém fluente em inglês).
\begin{resumo}[Abstract]
\begin{otherlanguage*}{english}
This study examines whether Brazilian listed family firms differ from non-family
firms in their use of debt. Drawing on the socioemotional wealth perspective, it
hypothesizes that controlling families, averse to creditor interference and to
the loss of control, adopt more conservative financing policies. The sample
comprises 16 healthcare companies listed on B3, with quarterly Economática data
from 2018 to 2025. Firms were classified as family firms following
\citeonline{santos2018}, and two alternative definitions were used as robustness
checks. Debt was measured by accounting ratios of leverage, debt service capacity
and debt maturity, and group differences were assessed with Welch's t test and the
Mann-Whitney test over the full period and over the pre-pandemic, pandemic and
post-pandemic subperiods. Family firms showed lower financial debt and a smaller
share of short-term debt in every year, but these differences were not
statistically significant over the full period, and total liabilities were
similar in both groups. The exception is debt maturity during the COVID-19
pandemic, when the smaller share of short-term debt in family firms was
significant at the 5\% level, in an exploratory subperiod analysis without
correction for multiple testing. The difference in financial debt changed sign
under the alternative family firm definitions and was sensitive to the exclusion
of one firm with no financial debt. The hypothesis is therefore not statistically
supported. The descriptive evidence is consistent with a more conservative stance
of families towards debt maturity, but also with differences in firm size and
access to credit between the groups.

\vspace{\onelineskip}
\noindent
\textbf{Keywords}: family firms; debt; capital structure; socioemotional wealth;
healthcare sector.
\end{otherlanguage*}
\end{resumo}
